\subsection{Small outside quotients and at least three columns}

\begin{lemma}[A small column image of the constant-action quotient]
\label{C:old-column-image-two-transfer}
Let $M$ have minimal band $I\times J$, and let $\bar S$ be
the constant-action quotient of a finite right act $S$.
Write $c_j\in\bar S$ for the constant value of an old band
element in column $j$, and put $T=\{c_j:j\in J\}$.
If $|T|=1$, every rational degree-two product of the original
act cone is zero.  If $|T|=2$, those products pull back from
an act cone over a monoid with minimal band $I\times T$ and
the same number of outside elements as $M$.
The latter comparison is an actual homotopy equivalence
after replacing $S$ by $\bar S$.
\end{lemma}
\begin{proof}
For $|T|=1$, the unique old column value $s_0$ is fixed by
every old actor.  Each transferred balanced cocycle has
\[
 u_L(f,g)=b^L_fa_g,\qquad
 b^L_m=L([s_0m]-[s_0],-)=0
\]
on old actors.  Its degree-two products therefore vanish
on all normalized act-category four-chains, as in
Theorem~\ref{C:three-by-three-act-cones21}.

For $|T|=2$, identify old columns with the same $c_j$ and
retain every outside element.  Equivariance of $j\mapsto c_j$
makes this a monoid quotient $M'$.  The act $\bar S$ factors
through it.  The compression monoids of
$(M,\bar S)$ and $(M',\bar S)$ are literally the same: both
have the old outside elements and the band $I\times\bar S$,
and every product and action is unchanged.
The compression theorem proves the asserted homotopy
comparison.  Finally
Lemma~\ref{C:constant-action-quotient-H2} is surjective on
degree-two cohomology in the pullback direction, which gives
the stated transfer to the original act cone.
\end{proof}

\begin{proposition}[Small outside quotients and small acts]
\label{C:three-outside-six-state-band-cones}
Let $M$ have rectangular-band minimal ideal $I\times J$ and
$q$ outside elements.  Its act cone cannot have type
$\mathbb CP^2$ in any of the following cases:
\[
 q=1;\qquad q\leq3\ \hbox{and }|S|\leq6;
 \qquad q\leq4\ \hbox{and }|S|\leq4.
\]
There is no restriction on the band dimensions.
\end{proposition}
\begin{proof}
For $q=1$, the contracted algebra is $\mathbb Q$, so the
row module is projective and
Lemma~\ref{C:row-pd-one-act-cups} applies at every state count.
Assume a projective-plane cone in one of the other cases.
It has no state fixed by the band, by the band-grid relative
argument.  Its act is connected, by the degree-one homology
sequence of the act-cone pushout.

If the band action is constant, compression gives an actual
projective-plane monoid with the same $q\leq4$ outside
elements.  The band Tor formula and
Theorem~\ref{B:contracted-four-Tor} exclude this.
If the action is not constant, there are at least two core
orbits.  Each has at least two states, each is individually
$M$-invariant, and their union cannot exhaust a connected
act.  Thus at most four states cannot occur at all.

For at most six states, at least one core orbit has exactly
two points: two orbits of three or more points would already
exhaust all six states.  The column values of that orbit
identify the old column image $T$ in the constant-action
quotient into at most two points.  Apply
Lemma~\ref{C:old-column-image-two-transfer}.
If $|T|=1$, the products are zero directly.  If $|T|=2$, the
comparison monoid has two columns and $q\leq3$ outside
elements, so
Lemma~\ref{C:three-outside-all-two-column-act-cups} makes
all its act-cone products zero without requiring bounded
homology for that intermediate cone.  Pullback gives the
same contradiction on the original projective plane.
\end{proof}

\begin{corollary}[At least three minimal columns below twenty-two arrows]
\label{C:three-or-more-column-act-cones21}
An act cone with rectangular-band minimal ideal having at
least three columns and at most twenty-one total arrows
cannot realize $\mathbb CP^2$.
\end{corollary}
\begin{proof}
At most two rows are excluded by the all-order two-row act
theorem, and at most three states are excluded below
twenty-two arrows.  Thus the band has at least three rows
and the act has at least four states.  The arrow budget and
the identity outside the band leave only
\[
 3\times3,\quad3\times4,\quad4\times3,
 \quad3\times5,\quad5\times3.
\]
The first shape is Theorem~\ref{C:three-by-three-act-cones21}.
For the last two, the budget forces $q=1$.
For the twelve-element bands, $1\leq q\leq4$ and
$|S|\leq8-q$.  The case $q=1$ is unrestricted; the cases
$q=2,3$ have at most six states; and $q=4$ has at most four.
The preceding proposition covers each case.
\end{proof}
